% chapters/discussion.tex
\chapter{Discussion}
\label{ch:discussion}

\section{Interpretation of Criterion Weights}
\label{sec:disc-weights}

The AHP weights show that graduate student reviewers clearly value trustworthiness
and completeness above the other criteria, together accounting for more than 75\% of
total weight. This ordering has a straightforward interpretation in an educational
context: an explanation is only useful as a learning tool if it is first credible and
adequately scoped. Reviewers stated in free-text comments that trustworthiness forms
the foundation for all other criteria — once an explanation fails the credibility
test, secondary qualities such as ease of following or goal support become irrelevant.

This pattern is consistent with the framework of \citet{hoffman2019metrics}, who
argue that trustworthiness is a prerequisite quality rather than one dimension among
equals. It also aligns with findings from social transparency research
\citep{ehsan2021expanding}, where users who encounter unreliable AI outputs reduce
engagement with all subsequent information from that system.

The subordinate position of usefulness and understandability does not imply that
these criteria are unimportant. Rather, the results indicate that they are activated
conditionally: high-performing models earned positive usefulness and understandability
ratings precisely because their explanations were first perceived as trustworthy.
Low-performing models received uniformly depressed scores across all criteria,
suggesting a ceiling effect driven by insufficient credibility.

\section{Model Performance and Practical Implications}
\label{sec:disc-models}

The interval-based ranking produces a stable ordering that differs in important ways
from what accuracy-only metrics would suggest. \texttt{qwen2.5:14b} is the strongest
model by a substantial margin. Its performance floor ($T_{Q1}$) far exceeds those of
all other models, and its IQR bands are narrow, indicating that it consistently
produces explanations that reviewers find reliable, complete, and useful across all
five reasoning task categories.

The finding that \texttt{phi3:3.8b} outperforms the larger \texttt{phi3:medium}
within the Phi family is notable. It suggests that architectural choices, including
training data composition, instruction-tuning strategy, and output formatting
conventions, have a stronger effect on explanation quality than raw parameter count.
This finding is relevant for practitioners choosing models under resource constraints:
a smaller model with appropriate architecture may outperform a larger but less
well-aligned one.

The fluency trap observed in \texttt{mistral:7b} is the most practically consequential
finding in this study. The gap between its Understandability score and its
Trustworthiness score is the largest of any model in the evaluation. For students
using this model's explanations to study economics, the risk is that coherent and
readable outputs reinforce incorrect causal reasoning rather than correcting it. This
is consistent with the concern raised by \citet{jin2024llm} that LLMs conflate
correlation with causation in superficially fluent ways, and by \citet{wu2024causality}
that reliability of causal claims is systematically difficult to assess from surface
form.

\section{Interval Width as Evidence of Reasoning Instability}
\label{sec:disc-interval}

One of the distinctive features of this study is the use of IQR intervals rather
than point estimates. The interval width carries substantive meaning: a wide interval
indicates that reviewers disagreed strongly about the quality of an explanation,
which in turn reflects ambiguity or inconsistency in the underlying model output.
This interpretation is supported by the correlation analysis, which finds
Trustworthiness and Usefulness as the dominant drivers of overall score
($r = 0.993$ and $r = 0.995$ respectively), consistent with the AHP weights.

Wide intervals are most common in Quantitative Reasoning and Procedural Knowledge
tasks, where the difference between a valid causal chain and a plausible narrative is
hardest to detect without domain expertise. Narrow intervals are most common in
Qualitative Reasoning, where the stylistic cues associated with correct reasoning are
more readily recognised by reviewers. This pattern suggests that the framework's
sensitivity to reviewer disagreement is itself informative about the nature of the
reasoning task.

\section{Implications for LLM Explanation Design}
\label{sec:disc-design}

The results have three direct implications for practitioners developing or deploying
LLMs in educational settings.

First, since trustworthiness received the highest weight, explanation systems should
make reliability signals explicit. This includes structured reasoning steps, explicit
statements of uncertainty, and transparent links between claims and supporting
evidence. The strongest model in this study (\texttt{qwen2.5:14b}) consistently
provided structured causal chains, which appears to drive both higher trustworthiness
and higher completeness ratings.

Second, model selection should be task-specific. No model excels uniformly across
all five reasoning task categories. Practitioners should identify the task types most
relevant to their application and select models accordingly, rather than relying on
aggregate benchmarks.

Third, fluency is not a reliable proxy for quality. Standard accuracy benchmarks do
not penalise models that produce coherent but unreliable explanations. The framework
introduced here provides a method for detecting this mismatch, which is important for
any deployment where explanation quality has direct consequences for student learning
or decision-making.

\section{Limitations}
\label{sec:disc-limitations}

Several limitations constrain the generalisability of these findings.

The reviewer pool consists of nine graduate students from a single institution with
a shared economics background. While this homogeneity increases within-study
consistency, it limits the generalisability of the derived criterion weights to other
reviewer populations, disciplines, or cultural contexts. Future work should replicate
the study with larger and more diverse reviewer pools.

The study focused on a single domain (Economics) and a single benchmark (MMLU).
Explanation preferences and criterion weights may differ in other academic disciplines,
particularly those with different conventions for causal argumentation.

Adaptive question selection increased the informativeness of the evaluation, but
prevents direct comparison with unselected questions. The 12 selected questions were
optimised for discriminative power and may not be representative of the full item
distribution.

Finally, all nine models evaluated are open-source. The framework has not been
validated on closed-source models, whose explanation characteristics may differ
substantially.
